% ============================================================================
% Fichier  : partie3/P03_C07_collecter_preuve.tex
% Titre    : Collecter la Preuve : OUI, NON, NEK
% Auteur   : MINKA MI NGUIDJOÏ Thierry Emmanuel
% Version  : 0.2 -- novembre 2026 (réécriture étoffée, alignée sur INF3226
%            C10, A0 et A1)
% ============================================================================
% Sources : A0 (hal-05701007), proposition de modalité d'ignorance explicite ;
% A1 (hal-05701270), instrument OUI/NON/NEK, mesures, indice rho ; GUID 5100 ;
% manuel IDI (documentation).
% Faits MINÉQUIR établis ici : domaine D1 Accès, 5 questions x 4 répondants
% (DAF EBOUMBOU, chef de la cellule ATANGANA, chef de la liquidation
% NDZIE, contrôleur financier MOUKOURI) ; grille binaire de l'audit interne
% 2025 : 17 OUI / 20 ; instrument trichotomique en mars 2026 : 11 OUI, 3 NON,
% 6 NEK -> mu 0,55, U 0,30, H 1,41 ; rho = 0,30.
% Faits CAPNO (LAB-05) : paie de décembre 2025, 4 312 lignes, net total
% 1 293 418 500 FCFA ; première extraction 4 309 lignes (3 manquantes,
% 1 905 000 FCFA) ; extraction refaite le 6 mai 2026.
% ============================================================================

\chapter[Collecter la preuve : OUI, NON, NEK]{Collecter la Preuve : OUI, NON, NEK}
\label{chap:preuve}
\chaptermark{Collecter la preuve}

\epigraph{%
    « Une grille qui interdit de répondre \emph{je ne sais pas}\\
    ne supprime pas l'ignorance~: elle la déguise en certitude. »%
}{Principe de l'auditologie}

\niveauChapitre{Intermédiaire}{Chapitre~\ref{chap:preparer}.}

\begin{boxobjectifs}
\begin{itemize}
    \item Utiliser l'instrument \OUI{}~/ \NON{}~/ \NEK{} auprès des
          répondants et pour coter les contrôles, et justifier la modalité
          \NEK{} à partir des axiomes de l'auditologie.
    \item Mesurer le biais d'une grille binaire par l'indice $\rho$.
    \item Classer les preuves selon leur fiabilité et pratiquer la
          corroboration.
    \item Conduire un entretien d'audit qui produit des réponses cotables.
    \item Calculer la taille d'un échantillon par attributs et en tirer une
          conclusion.
    \item Obtenir une extraction dont on garantit l'origine, l'exhaustivité,
          l'exactitude et l'intégrité.
    \item Tenir des papiers de travail qu'un tiers peut relire et refaire.
\end{itemize}
\end{boxobjectifs}

\bigskip

% ============================================================================
\section*{Pourquoi ce chapitre}
% ============================================================================

Dans le prologue, tout reposait sur deux fichiers et sur huit cotes. Si la
DAF avait pu dire «~ce fichier n'est pas le nôtre~», ou si la vérificatrice
avait coté \OUI{} ce qu'elle ne savait pas, le constat tombait. Ce chapitre
apprend à répondre à chaque question de la fiche d'assertion avec une
preuve qui tient, et à dire honnêtement ce que l'on ne sait pas. Il repose
sur l'instrument de réponse construit par l'article A1~\cite{minka2026nek},
à partir des axiomes de l'article A0~\cite{minka2026auditologie}.

\medskip
\textbf{Plan du chapitre.}
\begin{enumerate}
    \item L'instrument \OUI{}~/ \NON{}~/ \NEK{} (section~\ref{sec:c07-instrument}).
    \item Le biais du choix forcé (section~\ref{sec:c07-biais}).
    \item Les preuves et leur fiabilité (section~\ref{sec:c07-preuves}).
    \item Les entretiens (section~\ref{sec:c07-entretiens}).
    \item Tester et échantillonner (sections~\ref{sec:c07-techniques}
          et~\ref{sec:c07-echantillon}).
    \item Obtenir une extraction (section~\ref{sec:c07-extraction}).
    \item Les papiers de travail (section~\ref{sec:c07-papiers}).
\end{enumerate}

\separateur

% ============================================================================
\section{L'instrument OUI / NON / NEK}
\label{sec:c07-instrument}
% ============================================================================

\subsection{Pourquoi une troisième réponse}

L'article \emph{Vers l'auditologie} établit, à partir de quatre de ses
axiomes, que tout instrument de réponse d'audit valide doit comporter une
modalité explicite de connaissance insuffisante~\cite{minka2026auditologie}.
Le raisonnement est simple.
\begin{itemize}
    \item \textbf{Observabilité partielle} (axiome~1)~: l'auditeur ne voit
          qu'une partie des états du système~; certaines questions n'ont
          pas de réponse établie au moment où il les pose.
    \item \textbf{Cognition limitée} (axiome~2)~: les répondants, comme
          l'auditeur, ne savent pas tout de ce qu'ils gèrent.
    \item \textbf{Mesurabilité} (axiome~3)~: l'incertitude peut se mesurer,
          à condition d'être enregistrée.
    \item \textbf{Interaction stratégique} (axiome~4)~: un répondant a
          parfois intérêt à répondre \OUI{} quand il ne sait pas.
\end{itemize}
Une grille qui n'admet que \OUI{} et \NON{} oblige donc à transformer
l'ignorance en une réponse fausse. L'article A1 construit l'instrument qui
l'évite~\cite{minka2026nek}.

\begin{boxdef}
\termecle{Espace de réponse trichotomique}\index{espace de réponse
trichotomique}~: ensemble $\Omega = \{\OUI, \NON, \NEK\}$ des réponses
admises à une question d'audit.
\begin{itemize}
    \item \OUI{}~: le critère est satisfait, et une preuve l'établit.
    \item \NON{}~: le critère n'est pas satisfait, et une preuve l'établit.
    \item \NEK{} (\textit{Not Enough Knowledge})~: l'information disponible
          ne permet pas de conclure, malgré des diligences normales.
\end{itemize}
\end{boxdef}

\subsection{Deux usages}

\begin{enumerate}
    \item \textbf{Auprès des répondants}~: une même question est posée à
          plusieurs interlocuteurs. Pour chaque question, ou chaque domaine,
          on obtient une distribution de réponses, dont l'article A1 tire
          trois mesures~: $\mu$, la part des \OUI{}~; $U$, la part des
          \NEK{}~; $H$, l'entropie~\cite{minka2026nek}. Le
          Chapitre~\ref{chap:constater} apprend à les lire.
    \item \textbf{Par l'auditeur}~: après ses tests, l'auditeur cote chaque
          contrôle \OUI{}, \NON{} ou \NEK{}, preuve à l'appui. C'est la cote
          qui entre dans le rapport.
\end{enumerate}
Les deux usages se répondent~: une question où les répondants divergent
appelle des tests~; les tests permettent ensuite de coter avec preuve.

\subsection{NEK n'est pas « non applicable »}

\begin{boxalert}
\textbf{Violation typique.} Coter \NEK{} un contrôle qui ne s'applique pas au
périmètre (par exemple, la validation des annulations de recettes dans un
service qui n'encaisse rien).

\smallskip
\textbf{Règle~:} un contrôle sans objet est \emph{non applicable}~: il sort
de la décomposition, avec la raison écrite. \NEK{} est réservé à un contrôle
applicable dont on ne peut pas établir l'état. Les confondre fausse le taux
d'incertitude~$U$.
\end{boxalert}

\separateur

% ============================================================================
\section{Le biais du choix forcé}
\label{sec:c07-biais}
% ============================================================================

Une grille binaire force le répondant qui ne sait pas à choisir. Il choisit
le plus souvent la réponse attendue, ou la plus confortable~: le \OUI{}.
La grille binaire \emph{surestime} donc la conformité. L'article A1 mesure
cette surestimation par un indice de réduction du biais~\cite{minka2026nek}~:
\[
    \rho = \hat{\mu}^{B} - \mu^{T},
\]
où $\hat{\mu}^{B}$ est la part des \OUI{} obtenue avec l'instrument binaire, et
$\mu^{T}$ celle obtenue avec l'instrument trichotomique, sur les mêmes
questions et les mêmes répondants.

\begin{boxlecture}{Domaine D1, binaire puis trichotomique}
En 2025, l'audit interne du MINÉQUIR avait posé cinq questions sur les accès
à SIDÉPENSE à quatre interlocuteurs (le DAF, le chef de la cellule
informatique, le chef de la liquidation, le contrôleur financier), avec une
grille \OUI{}~/ \NON{}~: 17 \OUI{} sur 20 réponses. En mars~2026, la
vérificatrice pose les mêmes questions aux mêmes personnes avec
l'instrument trichotomique~: 11 \OUI{}, 3 \NON{}, 6 \NEK{}.

\medskip
\textbf{Constatation.} $\hat{\mu}^{B} = 17/20 = 0{,}85$~; $\mu^{T} = 11/20 =
0{,}55$~; $\rho = 0{,}30$.

\medskip
\textbf{Analyse} (confiance moyenne). La grille binaire surestimait la
conformité de trente points~: six réponses \OUI{} de 2025 étaient en
réalité des «~je ne sais pas~». C'est ainsi que l'audit interne avait pu
écrire que «~la gestion des accès est satisfaisante~», quelques mois avant
le double paiement.

\medskip
\hfill\lecture{0{,}55}{0{,}30}{1{,}41}
\end{boxlecture}

\begin{boxinfo}
\textbf{Ce que $\rho$ ne dit pas.} L'indice mesure l'écart entre deux
instruments, pas la conformité réelle. Seuls les tests de l'auditeur disent
si les contrôles fonctionnent~: au Chapitre~\ref{chap:acces}, un seul des
six contrôles d'accès sera coté \OUI{}.
\end{boxinfo}

\separateur

% ============================================================================
\section{Les preuves et leur fiabilité}
\label{sec:c07-preuves}
% ============================================================================

\subsection{Suffisante et appropriée}

\begin{boxdef}
\termecle{Preuve d'audit}\index{preuve d'audit}~: information sur laquelle
l'auditeur fonde sa cote et ses conclusions. Elle doit être
\emph{suffisante} (en quantité, au regard du risque) et \emph{appropriée}
(pertinente pour le critère, et fiable).
\end{boxdef}

\subsection{Une hiérarchie de fiabilité}

\begin{tableaucours}{cL{4.2cm}L{4.4cm}L{2.8cm}}{%
    \tableauHeader{Rang} &
    \tableauHeader{Source} &
    \tableauHeader{Exemple} &
    \tableauHeader{Fiabilité}
}
1 & Trace du système, extraite par l'auditeur & Matrice des profils
extraite devant la vérificatrice & La plus élevée \\
2 & Document externe & Relevé bancaire, visa du poste comptable & Élevée \\
3 & Document interne approuvé & Note de service signée & Moyenne \\
4 & Observation directe & La vérificatrice voit un agent liquider &
Ponctuelle \\
5 & Déclaration orale & « Les comptes sont désactivés dès qu'on nous
informe » & La plus faible \\
\end{tableaucours}
\vspace{-0.8em}
\begin{center}{\small\itshape Tableau~7.1. Fiabilité des preuves selon leur
source}\end{center}

Une déclaration orale ne suffit jamais à coter \OUI{}~: elle oriente la
recherche de la preuve. Quand l'auditeur ne dispose que d'une déclaration,
la cote est \NEK{}, et la question devient «~où est la preuve~?~».

\subsection{La corroboration}

Deux preuves indépendantes qui concordent valent plus que la plus fiable
des deux prise seule. Pour les contrôles importants, l'auditeur cherche une
seconde source~: la liste des mutations de la DRH (rang~3) recoupée avec
les dates de désactivation dans SIDÉPENSE (rang~1). Si deux sources se
contredisent sans qu'on puisse les départager, la cote est \NEK{}, et
l'écart devient lui-même un objet d'investigation.

\separateur

% ============================================================================
\section{Les entretiens}
\label{sec:c07-entretiens}
% ============================================================================

L'entretien d'audit, à la différence de l'entretien de prise de
connaissance (Chapitre~\ref{chap:preparer}), recueille des réponses
cotables. Ses questions sont celles des fiches d'assertion, classées par
interlocuteur.

\begin{boxprocedure}
\textbf{Règles de l'entretien d'audit}~:
\begin{enumerate}
    \item un répondant à la fois, pour éviter qu'un chef ne réponde pour son
          équipe~;
    \item lire la question de la fiche telle quelle, sans la reformuler en
          suggérant la réponse~;
    \item annoncer que \NEK{} est une réponse admise et utile~;
    \item noter la réponse, son auteur, la date, et la pièce invoquée~;
    \item demander la pièce séance tenante, ou fixer un délai écrit~;
    \item faire relire et signer le compte rendu.
\end{enumerate}
\end{boxprocedure}

\begin{boxalert}
\textbf{Violation typique.} Demander au DAF «~vous désactivez bien les
comptes des agents mutés, n'est-ce pas~?~».

\smallskip
\textbf{Règle~:} une question qui suggère la réponse produit un \OUI{} sans
valeur. On demande «~qui désactive, quand, sur quelle pièce~?~», et l'on
admet «~je ne sais pas~».
\end{boxalert}

\separateur

% ============================================================================
\section{Les techniques de test}
\label{sec:c07-techniques}
% ============================================================================

\begin{tableaucours}{L{2.6cm}L{4.4cm}L{5.2cm}}{%
    \tableauHeader{Technique} &
    \tableauHeader{Ce qu'elle établit} &
    \tableauHeader{Exemple}
}
Entretien & Ce que l'entité déclare & Le chef de la cellule décrit la
gestion des comptes \\
Observation & Ce qui se passe, au moment où l'on regarde & Assister à une
désactivation de compte \\
Inspection & Ce que contient un document ou un enregistrement & Lire la
matrice des profils \\
Ré-exécution & Que le contrôle produit le résultat attendu & Tenter de
saisir un mandat en double dans un environnement de test \\
Analyse de données & Ce que dit toute la population & Rechercher les
doublons parmi 9\,611 mandats \\
\end{tableaucours}
\vspace{-0.8em}
\begin{center}{\small\itshape Tableau~7.2. Les techniques de
test}\end{center}

\begin{boxalert}
\textbf{Violation typique.} Faire une ré-exécution dans l'application de
production («~pour voir si le doublon passe~»).

\smallskip
\textbf{Règle~:} l'auditeur ne crée, ne modifie ni ne supprime jamais une
donnée de production. Les essais se font dans un environnement de test, ou
par observation d'une opération réelle.
\end{boxalert}

\separateur

% ============================================================================
\section{Échantillonner dans le système}
\label{sec:c07-echantillon}
% ============================================================================

\subsection{Population et échantillon}

\begin{boxdef}
\termecle{Population}\index{population}~: ensemble des opérations sur
lesquelles porte la conclusion. Dans un SI, elle se prend dans le système,
après rapprochement de ses totaux (Chapitre~\ref{chap:prologue}).

\medskip
\termecle{Échantillon}\index{échantillon}~: partie de la population
examinée en détail. Il est \emph{aléatoire} quand chaque opération a la même
chance d'être tirée, \emph{raisonné} quand l'auditeur choisit les opérations
selon un critère (montant élevé, opération hors heures).
\end{boxdef}

L'échantillon aléatoire permet d'étendre la conclusion à la population~;
l'échantillon raisonné va droit aux anomalies probables, mais ne permet pas
cette extension. L'auditeur dit toujours lequel il a utilisé.

\subsection{L'échantillonnage par attributs, démontré}

On teste un attribut (le contrôle a-t-il été exécuté~?) sur $n$ opérations
tirées au hasard. On veut pouvoir dire, avec une confiance $c$, que le taux
réel d'écart ne dépasse pas un taux toléré $t$, si l'on ne trouve aucun
écart.

Si le taux réel d'écart était exactement $t$, la probabilité de ne trouver
aucun écart sur $n$ tirages indépendants serait $(1 - t)^n$. Pour que ce
résultat soit peu probable, au plus $1 - c$, il faut~:
\[
    (1 - t)^n \leq 1 - c
    \quad\Longleftrightarrow\quad
    n \geq \frac{\ln(1 - c)}{\ln(1 - t)}.
\]

\begin{tableaucours}{L{3.0cm}L{3.0cm}L{3.0cm}L{3.0cm}}{%
    \tableauHeader{Confiance $c$} &
    \tableauHeader{Taux toléré $t$} &
    \tableauHeader{Calcul} &
    \tableauHeader{Taille $n$}
}
95~\% & 5~\% & $\ln 0{,}05 / \ln 0{,}95 \approx 58{,}4$ & 59 \\
90~\% & 5~\% & $\ln 0{,}10 / \ln 0{,}95 \approx 44{,}9$ & 45 \\
95~\% & 10~\% & $\ln 0{,}05 / \ln 0{,}90 \approx 28{,}4$ & 29 \\
\end{tableaucours}
\vspace{-0.8em}
\begin{center}{\small\itshape Tableau~7.3. Tailles d'échantillon par
attributs, sans écart attendu}\end{center}

\subsection{Conclure sur l'échantillon}

\begin{itemize}
    \item \textbf{Aucun écart}~: avec une confiance $c$, le taux d'écart ne
          dépasse pas $t$~; le contrôle est coté \OUI{} pour la période.
    \item \textbf{Un écart ou plus}~: la conclusion par attributs tombe~; le
          contrôle est \NON{}, et l'on examine chaque écart (cause, montant).
    \item \textbf{Des éléments introuvables}~: s'ils sont nombreux, la cote
          est \NEK{}, et l'on écrit les diligences faites.
\end{itemize}

\begin{boxoutils}{Tirer un échantillon aléatoire dans un tableur}
Dans la feuille des mandats, ajouter une colonne~:
\begin{lstlisting}
H2 : =ALEA()
\end{lstlisting}
Copier la colonne et la coller en valeurs (pour figer le tirage), trier sur
H, retenir les $n$ premières lignes. Conserver la feuille triée~: elle
prouve le caractère aléatoire du tirage.
\end{boxoutils}

\subsection{Tester toute la population}

Quand la donnée est fiable, il est souvent plus probant de tester
\emph{toute} la population par une analyse~: doublons, trous de
numérotation, opérations hors heures, montants juste sous un seuil
(Chapitre~\ref{chap:donnees-fraude}). L'analyse désigne des opérations~;
chacune est ensuite examinée pièce par pièce.

\separateur

% ============================================================================
\section{Obtenir une extraction}
\label{sec:c07-extraction}
% ============================================================================

\subsection{Quatre garanties}

\begin{tableaucours}{L{2.8cm}L{4.4cm}L{5.0cm}}{%
    \tableauHeader{Garantie} &
    \tableauHeader{Question} &
    \tableauHeader{Comment l'obtenir}
}
Origine & Le fichier vient-il bien de l'application~? & Assister à
l'extraction, ou la faire faire sous ses yeux \\
Exhaustivité & Toutes les opérations y sont-elles~? & Requête écrite~;
nombre de lignes et totaux rapprochés d'une source indépendante \\
Exactitude & Les valeurs sont-elles celles de l'application~? & Comparer
quelques lignes à l'écran \\
Intégrité & Le fichier n'a-t-il pas changé depuis~? & Empreinte SHA-256
calculée à la remise, contresignée \\
\end{tableaucours}
\vspace{-0.8em}
\begin{center}{\small\itshape Tableau~7.4. Les quatre garanties d'une
extraction}\end{center}

\subsection{L'empreinte et les totaux de contrôle}

\begin{boxdef}
\termecle{Empreinte}\index{empreinte} (\textit{hash})~: suite de caractères
calculée à partir du contenu d'un fichier par une fonction de hachage
(SHA-256). Le moindre changement du fichier change l'empreinte.

\medskip
\termecle{Totaux de contrôle}\index{totaux de contrôle}~: nombre de lignes
et sommes, rapprochés d'une source que l'auditeur n'a pas obtenue de la même
personne (relevé bancaire, compte administratif, état du poste comptable).
\end{boxdef}

L'empreinte prouve que le fichier n'a pas changé~; elle ne prouve pas qu'il
était complet. Il faut les deux.

\begin{boxoutils}{Calculer et vérifier une empreinte}
Sous Windows, dans l'invite de commandes~:
\begin{lstlisting}
certutil -hashfile paie_2025-12.csv SHA256
\end{lstlisting}
Sous Linux (bureau d'analyse de NDAP SUKULU)~:
\begin{lstlisting}
sha256sum paie_2025-12.csv
sha256sum -c MANIFESTE_SHA256.txt
\end{lstlisting}
\end{boxoutils}

\begin{boxscenario}{Dossier CAPNO. Trois lignes de moins}
\textbf{Constatation.} La première extraction de la paie de décembre~2025,
remise sur une clé, compte 4\,309 lignes. L'ordre de virement transmis à la
banque porte sur 1\,293\,418\,500~FCFA pour 4\,312 bénéficiaires. Il manque
trois lignes, pour un total de 1\,905\,000~FCFA.

\medskip
\textbf{Analyse} (confiance élevée). L'extraction n'était pas exhaustive.
M.~HAMADOU la fait refaire devant lui le 6~mai~2026~: 4\,312 lignes, totaux
égaux à l'ordre de virement. Les trois lignes absentes deviennent les
premières à examiner~: pourquoi manquaient-elles~?
\end{boxscenario}

\separateur

% ============================================================================
\section{Les papiers de travail}
\label{sec:c07-papiers}
% ============================================================================

Le manuel de l'IDI le rappelle~: la documentation doit permettre à un
auditeur expérimenté, sans lien avec la mission, de comprendre les travaux,
les preuves et les conclusions~\cite{idi2021conformite}. C'est aussi la
traduction pratique de l'axiome d'\emph{opposabilité}~: une conclusion qui
ne peut pas être refaite ne peut pas être défendue~\cite{minka2026auditologie}.

\begin{boxprocedure}
\textbf{Feuille de test}, une par contrôle testé~:
\begin{enumerate}
    \item référence (position dans le gabarit, ligne du programme)~;
    \item fiche d'assertion (critères et questions)~;
    \item population (source, nombre, totaux, empreinte du fichier)~;
    \item échantillon (méthode, confiance, taux toléré, taille, liste)~;
    \item résultat élément par élément, avec la pièce examinée~;
    \item cote de chaque critère et cote du contrôle, avec la
          constatation~; pour un \NEK{}, les diligences et la relance~;
    \item auteur, date, et visa du chef de mission.
\end{enumerate}
\end{boxprocedure}

\begin{boxinfo}
\textbf{Conserver.} Les extractions sont des données personnelles (paie) ou
sensibles (fournisseurs). Elles sont conservées sur un support protégé, avec
leur empreinte, aussi longtemps que les suites possibles l'exigent, puis
détruites selon les règles de l'institution.
\end{boxinfo}

\begin{boxmetier}{Vérificateur}
Annoncez \NEK{} dès le premier entretien~: «~si vous ne savez pas,
dites-le~; c'est une réponse utile~». Vous obtiendrez moins de \OUI{}, et
davantage de vérité.
\end{boxmetier}

\begin{boxmetier}{Juge des comptes}
Une pièce électronique produite à l'instance vaut si l'on peut en établir
l'origine et l'intégrité. Demandez le procès-verbal de remise et
l'empreinte~; leur absence affaiblit la pièce.
\end{boxmetier}

\begin{boxmetier}{Auditeur interne}
Remplacez vos grilles \OUI{}~/ \NON{} par l'instrument à trois réponses, et
mesurez $\rho$ sur vos propres questionnaires~: vous saurez de combien vos
évaluations passées surestimaient la conformité.
\end{boxmetier}

\separateur

\begin{boxresume}
\begin{itemize}
    \item \textbf{Trois réponses, fondées en théorie.} \OUI{}, \NON{}, \NEK{}~;
          la modalité \NEK{} découle des axiomes~1 à~4 d'A0.
    \item \textbf{NEK n'est pas « non applicable ».} Un contrôle sans objet
          sort de la décomposition.
    \item \textbf{La grille binaire surestime}~; $\rho$ mesure de combien
          (0,30 sur les accès du MINÉQUIR).
    \item \textbf{Une hiérarchie de preuves}~; une déclaration seule donne
          \NEK{}~; deux sources indépendantes valent mieux qu'une.
    \item \textbf{L'entretien ne suggère pas}~; il admet «~je ne sais pas~».
    \item \textbf{$n \ge \ln(1-c) / \ln(1-t)$}~: 59 tirages pour 95~\% de
          confiance et 5~\% d'écart toléré.
    \item \textbf{Empreinte et totaux}~: intégrité et exhaustivité d'une
          extraction.
    \item \textbf{Un tiers doit pouvoir refaire}~: c'est l'opposabilité.
\end{itemize}
\end{boxresume}

\begin{boxexo}[title={\ding{45}\quad Exercice~7.1. Coter avec la bonne preuve \niveaubase}]
Pour chaque situation, donnez la cote et le rang de la preuve
(Tableau~7.1)~: (a) le DAF affirme que les profils sont revus chaque
semestre~; (b) la matrice extraite devant vous montre cinq comptes LIQ +
ORD~; (c) une note signée prévoit la désactivation sous 48~heures, mais le
journal n'a pas encore été remis~; (d) la DRH et la cellule donnent deux
dates de désactivation différentes pour le même agent.
\end{boxexo}

\begin{boxexo}[title={\ding{45}\quad Exercice~7.2. L'indice $\rho$ \niveaubase}]
Sur le domaine D4 (exploitation), la grille binaire de 2025 donnait 14
\OUI{} sur 20~; l'instrument trichotomique donne 8 \OUI{}, 4 \NON{}, 8
\NEK{}. Calculez $\hat{\mu}^{B}$, $\mu^{T}$ et $\rho$, et dites ce que cela
signifie.
\end{boxexo}

\begin{boxexo}[title={\ding{45}\quad Exercice~7.3. Dimensionner \niveaumoyen}]
Vous voulez conclure, avec 95~\% de confiance, que le visa du contrôle
financier manque sur moins de 3~\% des mandats de 2025. (a) Combien de
mandats tirez-vous~? (b) Vous trouvez un mandat sans visa~: que
concluez-vous~? (c) Que proposez-vous plutôt, puisque l'extraction des
9\,611 mandats est disponible~?
\end{boxexo}

\begin{boxexo}[title={\ding{45}\quad Exercice~7.4. Trois lignes \niveauexpert}]
Les trois lignes absentes de la première extraction de la CAPNO
correspondent à trois agents créés le 28~novembre~2025 par le même compte.
Dites ce que cela suggère, ce que vous ne pouvez pas encore affirmer, les
trois pièces que vous demandez, et la cote provisoire du contrôle
«~l'extraction remise est exhaustive~».
\end{boxexo}

\begin{boxinfo}
\textbf{Transition.} Les réponses sont recueillies, les contrôles cotés,
les preuves conservées. Le Chapitre~\ref{chap:constater} apprend à lire
ces réponses par trois mesures, à conclure, à rapporter et à suivre.
\end{boxinfo}
